% Source: https://artofproblemsolving.com/community/c3h1558984_easy_endless
% Stable identity: aops-topic-1558984-q05
Suppose $f$ has integer coefficients and $f(n)$ is prime for every positive integer $n$. Let $p=f(1)$, a prime.

For integers $u$, $v$, the difference $u-v$ divides $f(u)-f(v)$ (each term $c\left(u^j-v^j\right)$ is divisible by $u-v$). Taking $u=1+kp$ and $v=1$, for every positive integer $k$:
\[
p\mid f(1+kp)-f(1)=f(1+kp)-p ,\qquad\text{so}\qquad p\mid f(1+kp).
\]
But $f(1+kp)$ is prime, so it equals $p$. Thus the polynomial $f(x)-p$ has the infinitely many roots $1+p,1+2p,1+3p,\dots$, so it is the zero polynomial and $f$ is constant. This contradicts the assumption that $f$ is non-constant. $\blacksquare$

(Polynomials can produce many primes for a while: $n^2+n+41$ is prime for $n=0,1,\dots,39$, but $40^2+40+41=41^2$.)
